\documentclass[11pt,a4paper]{article}
\usepackage[T1]{fontenc}
\usepackage[utf8]{inputenc}
\usepackage{lmodern}
\usepackage[margin=22mm]{geometry}
\usepackage{amsmath,booktabs,longtable}
\usepackage{xcolor,listings,textcomp}
\usepackage[hidelinks]{hyperref}
\hypersetup{pdftitle={Tutorial 2: MV Effects and Reference Voltage - Exercise Solutions},pdfauthor={},pdfsubject={Exercise Solutions}}
\definecolor{codeblue}{RGB}{29,72,117}
\definecolor{codegreen}{RGB}{40,105,68}
\definecolor{codebackground}{RGB}{247,248,250}
\lstset{
  language=Python,
  basicstyle=\ttfamily\footnotesize,
  keywordstyle=\color{codeblue}\bfseries,
  commentstyle=\color{codegreen},
  stringstyle=\color{codeblue},
  backgroundcolor=\color{codebackground},
  frame=single,rulecolor=\color{black!15},
  breaklines=true,breakatwhitespace=false,
  columns=fullflexible,keepspaces=true,
  showstringspaces=false,upquote=true,tabsize=4,
  xleftmargin=4pt,xrightmargin=4pt,
  aboveskip=10pt,belowskip=10pt
}
\setlength{\parindent}{0pt}
\setlength{\parskip}{6pt}
\title{Tutorial 2: MV Effects and Reference Voltage\\Exercise Solutions}
\author{}
\date{}

\begin{document}
\maketitle

\section*{Before running the solutions}
Open \path{Tutorial_2_MV_Effects_Reference_Voltage.ipynb} and, in a fresh
Python kernel, run Sections 2.2--2.4. These sections load the arrays and
define the model, scaling, validation and search functions used below.
There is no need to run the tutorial's Sections 2.5--2.6 first.
Paste the seven Python listings into the matching Exercise Workspace cells
and run them in order. The listings are complete code, not pseudocode.

These solutions follow the current notebook's data convention: all 31 columns
of \path{V.pkl} are targets; the three columns of \path{V_secondary.pkl}
are additional supplied reference inputs. Thus the two models have 62 and
65 inputs, respectively, and the same 31 outputs. The precise physical
extraction node of the secondary-reference file is not documented in the
supplied data package. No customer column is removed without a verified
reference-to-customer mapping.

Both exercises use ReLU, 40 epochs, three blocked folds, a learning rate of
$10^{-3}$, L2 strength $10^{-5}$ and batch size 72. This exercise budget is
independent of the tutorial's quick/full mode. All selection uses training
data only. Test data are evaluated only after choosing the configuration
and seed. This is a classroom reuse of the tutorial's test period, not a
new independent test set.

\section*{E1.1: Prepare the two input sets and common targets}
Keep the P/Q columns from all 31 customers in both input sets. Append the
three supplied reference columns only to the augmented model. Do not fit
scalers here: validation must fit them separately inside each training fold.

\begin{lstlisting}[title={E1.1}]
e1_reference_columns = REFERENCE_PHASES.copy()
e1_target_indices = np.arange(C, dtype=int)
e1_target_customers = e1_target_indices + 1
e1_model_inputs = {
    "PQ_only": PQ_mv_tr.astype(np.float32, copy=False),
    "PQ_plus_Vref": np.column_stack([
        PQ_mv_tr, V_secondary_mv_tr
    ]).astype(np.float32),
}
e1_targets = V_mv_tr[:, e1_target_indices].astype(
    np.float32, copy=False
)
e1_model_names = list(e1_model_inputs)

assert e1_reference_columns == [
    "phase_a", "phase_b", "phase_c"
]
assert e1_model_inputs["PQ_only"].shape == (6048, 62)
assert e1_model_inputs["PQ_plus_Vref"].shape == (6048, 65)
assert e1_targets.shape == (6048, 31)
assert np.array_equal(
    e1_model_inputs["PQ_plus_Vref"][:, 62:],
    V_secondary_mv_tr,
)

display(pd.DataFrame([
    {
        "Model": MODEL_LABELS[name],
        "Training rows": len(e1_targets),
        "Inputs": e1_model_inputs[name].shape[1],
        "Outputs": e1_targets.shape[1],
    }
    for name in e1_model_names
]))
print("Reference columns:", e1_reference_columns)
print("Target customers:", e1_target_customers.tolist())

# Keep exercise results separate from tutorial results.
e1_results_directory = (
    TUTORIAL_DIRECTORY / "results" / "exercises"
)
e1_results_directory.mkdir(parents=True, exist_ok=True)
\end{lstlisting}

\section*{E1.2: Compare 93 and 155 hidden neurons}
Use the existing \texttt{run\_grid\_search()} function. It builds fresh
models and fits scalers on each fold's training rows only. The two
configurations require $2\times3\times2=12$ model fits.

\begin{lstlisting}[title={E1.2}]
e1_epochs = 40
e1_seed = 42
e1_k = 3
e1_search_space = {
    "hidden_neurons": [93, 155],
    "activation": ["relu"],
    "lr": [1e-3],
    "weight_decay": [1e-5],
    "batch_size": [72],
    "epochs": [e1_epochs],
}
e1_results = run_grid_search(
    e1_model_inputs,
    e1_targets,
    e1_search_space,
    k=e1_k,
    seed=e1_seed,
)
e1_results.to_csv(
    e1_results_directory / "exercise_configurations.csv",
    index=False,
)
\end{lstlisting}

\section*{E1.3: Display validation scores and select a configuration}
The joint score is the arithmetic mean of the two models' mean validation
RMSE values. Select one shared configuration using its unrounded score.
The reported standard deviations describe the three fold scores
(\texttt{numpy.std}, \texttt{ddof=0}), not confidence intervals.

\begin{lstlisting}[title={E1.3}]
e1_results = e1_results.sort_values(
    "joint_RMSE_kfold", kind="stable"
).reset_index(drop=True)
e1_score_columns = [
    "hidden_neurons", "activation", "epochs",
    "PQ_only_RMSE_kfold", "PQ_only_RMSE_std",
    "PQ_plus_Vref_RMSE_kfold", "PQ_plus_Vref_RMSE_std",
    "joint_RMSE_kfold",
]
display(e1_results[e1_score_columns].round(6))

e1_best_row = e1_results.iloc[0]
e1_best_config = {
    "hidden_neurons": int(e1_best_row["hidden_neurons"]),
    "activation": str(e1_best_row["activation"]),
    "lr": float(e1_best_row["lr"]),
    "weight_decay": float(e1_best_row["weight_decay"]),
    "batch_size": int(e1_best_row["batch_size"]),
    "epochs": int(e1_best_row["epochs"]),
}
assert e1_best_config["activation"] == "relu"
assert e1_best_config["epochs"] == 40
print("Selected exercise configuration:", e1_best_config)
\end{lstlisting}

\section*{E2.1: Compare seeds 0, 1 and 2}
Keep the configuration and folds fixed. This step uses training-period
data only and requires $3\times3\times2=18$ additional model fits.

\begin{lstlisting}[title={E2.1}]
e2_seeds = [0, 1, 2]
e2_seed_results = run_seed_search(
    e1_model_inputs,
    e1_targets,
    e1_best_config,
    e2_seeds,
    k=e1_k,
)
e2_seed_results = e2_seed_results.sort_values(
    "joint_RMSE_kfold", kind="stable"
).reset_index(drop=True)
display(e2_seed_results[[
    "seed",
    "PQ_only_RMSE_kfold", "PQ_only_RMSE_std",
    "PQ_plus_Vref_RMSE_kfold", "PQ_plus_Vref_RMSE_std",
    "joint_RMSE_kfold",
]].round(6))
e2_selected_seed = int(e2_seed_results.iloc[0]["seed"])
print("Selected exercise seed:", e2_selected_seed)
e2_seed_results.to_csv(
    e1_results_directory / "exercise_seeds.csv",
    index=False,
)
\end{lstlisting}

\section*{E2.2: Fit new scalers and train two new final models}
Fit one common target scaler and one input scaler per model on all exercise
training rows. Do not reuse the tutorial's final models or fold models.
These two fits bring the exercise total to $12+18+2=32$.

\begin{lstlisting}[title={E2.2}]
e2_output_scaler = MinMaxScaler().fit(e1_targets)
e2_y_train_scaled = e2_output_scaler.transform(
    e1_targets
).astype(np.float32)
e2_models = {}
e2_input_scalers = {}
e2_histories = {}
e2_training_rows = []

for name in e1_model_names:
    e2_input_scalers[name] = MinMaxScaler().fit(
        e1_model_inputs[name]
    )
    e2_x_train_scaled = e2_input_scalers[name].transform(
        e1_model_inputs[name]
    ).astype(np.float32)
    e2_models[name] = build_model(
        e2_x_train_scaled.shape[1],
        e2_y_train_scaled.shape[1],
        e1_best_config,
        seed=e2_selected_seed,
    )
    e2_histories[name] = train_model(
        e2_models[name],
        e2_x_train_scaled,
        e2_y_train_scaled,
        e1_best_config,
    )
    e2_completed_epochs = len(e2_histories[name].epoch)
    assert e2_completed_epochs == e1_best_config["epochs"]
    e2_training_rows.append({
        "Model": MODEL_LABELS[name],
        "Inputs": e2_models[name].input_shape[-1],
        "Outputs": e2_models[name].output_shape[-1],
        "Completed epochs": e2_completed_epochs,
    })

display(pd.DataFrame(e2_training_rows))
e2_run_settings = {
    "configuration": e1_best_config,
    "selected_seed": e2_selected_seed,
    "folds": e1_k,
    "training_fits": 32,
    "reference_source": "V_secondary.pkl",
    "reference_columns": e1_reference_columns,
    "target_customers": e1_target_customers.tolist(),
}
with (e1_results_directory / "exercise_settings.json").open(
    "w", encoding="utf-8"
) as settings_file:
    json.dump(e2_run_settings, settings_file, indent=2)
\end{lstlisting}

\section*{E2.3: Predict test voltages and report errors}
Transform test inputs with the fitted training scalers. Never call
\texttt{fit()} or \texttt{fit\_transform()} on test data.
The predictions are converted back to volts before calculating errors.
For $N$ time steps and $C$ targets, the overall metrics are
\[
  \mathrm{RMSE}=\sqrt{\frac{1}{NC}\sum_{t=1}^{N}\sum_{c=1}^{C}
  (\widehat V_{t,c}-V_{t,c})^2},\qquad
  \mathrm{MAE}=\frac{1}{NC}\sum_{t=1}^{N}\sum_{c=1}^{C}
  |\widehat V_{t,c}-V_{t,c}|.
\]
Per-customer RMSE averages over time only.

\begin{lstlisting}[title={E2.3}]
e2_test_inputs = {
    "PQ_only": PQ_mv_te.astype(np.float32, copy=False),
    "PQ_plus_Vref": np.column_stack([
        PQ_mv_te, V_secondary_mv_te
    ]).astype(np.float32),
}
e2_test_targets = V_mv_te[:, e1_target_indices].astype(
    np.float32, copy=False
)
assert e2_test_inputs["PQ_only"].shape == (2016, 62)
assert e2_test_inputs["PQ_plus_Vref"].shape == (2016, 65)
assert e2_test_targets.shape == (2016, 31)

e2_predictions = {}
e2_errors = {}
e2_metric_rows = []
for name in e1_model_names:
    e2_x_test_scaled = e2_input_scalers[name].transform(
        e2_test_inputs[name]
    ).astype(np.float32)
    e2_predictions[name] = predict_unscaled(
        e2_models[name], e2_x_test_scaled, e2_output_scaler
    )
    assert e2_predictions[name].shape == e2_test_targets.shape
    assert np.isfinite(e2_predictions[name]).all()
    e2_errors[name] = e2_predictions[name] - e2_test_targets
    e2_metric_rows.append({
        "Model": MODEL_LABELS[name],
        "RMSE (V)": float(np.sqrt(np.mean(
            e2_errors[name] ** 2
        ))),
        "MAE (V)": float(np.mean(np.abs(e2_errors[name]))),
    })

e2_metrics = pd.DataFrame(e2_metric_rows)
display(e2_metrics.round(6))
e2_customer_rmse = pd.DataFrame({
    "Customer": e1_target_customers,
    "P/Q only RMSE (V)": np.sqrt(np.mean(
        e2_errors["PQ_only"] ** 2, axis=0
    )),
    "P/Q + Vref RMSE (V)": np.sqrt(np.mean(
        e2_errors["PQ_plus_Vref"] ** 2, axis=0
    )),
})
display(e2_customer_rmse.round(6))
e2_metrics.to_csv(
    e1_results_directory / "exercise_test_metrics.csv",
    index=False,
)
e2_customer_rmse.to_csv(
    e1_results_directory / "exercise_customer_rmse.csv",
    index=False,
)
\end{lstlisting}

\section*{E2.4: Inspect the worst P/Q-only customer}
Choose the customer using the P/Q-only test RMSE, then use that same
customer for both model comparisons. This is post-training inspection,
not a further model-selection step. Residuals are estimate minus reference.

\begin{lstlisting}[title={E2.4}]
e2_worst_position = int(np.argmax(
    e2_customer_rmse["P/Q only RMSE (V)"].to_numpy()
))
e2_worst_customer = int(
    e1_target_customers[e2_worst_position]
)
print("Worst P/Q-only customer:", e2_worst_customer)
display(e2_customer_rmse.iloc[[e2_worst_position]].round(6))

e2_test_hours = (
    np.arange(len(e2_test_targets)) * MINUTES_PER_STEP / 60
)
e2_fig, e2_axes = plt.subplots(
    2, 1, figsize=(11, 6), sharex=True
)
e2_axes[0].plot(
    e2_test_hours,
    e2_test_targets[:, e2_worst_position],
    color="black", label="Simulated reference",
)
for name, colour in [
    ("PQ_only", "tab:red"),
    ("PQ_plus_Vref", "tab:green"),
]:
    e2_axes[0].plot(
        e2_test_hours,
        e2_predictions[name][:, e2_worst_position],
        color=colour, label=MODEL_LABELS[name], alpha=0.8,
    )
    e2_axes[1].plot(
        e2_test_hours,
        e2_errors[name][:, e2_worst_position],
        color=colour, label=MODEL_LABELS[name], linewidth=0.8,
    )
e2_axes[0].set_ylabel("Voltage (V)")
e2_axes[0].set_title(
    f"Exercise: worst P/Q-only customer {e2_worst_customer}"
)
e2_axes[0].legend()
e2_axes[1].axhline(0, color="black", linestyle="--")
e2_axes[1].set_xlabel("Hours from start of test period")
e2_axes[1].set_ylabel("Estimate - reference (V)")
e2_axes[1].set_title("Voltage residuals")
e2_axes[1].legend()
e2_fig.tight_layout()
e2_fig.savefig(
    e1_results_directory / "exercise_worst_customer.png",
    dpi=160, bbox_inches="tight",
)
plt.show()
\end{lstlisting}

\section*{Verified execution example}
The following results are from an actual run of the listings above, not
the tutorial's tanh models. Values may vary slightly across environments.
The code selects settings from validation scores; the numbers below must
not be hard-coded as a replacement for that selection.

All 32 exercise fits completed without errors.

\subsection*{E1.2--E1.3: Configuration comparison}

Mean and population standard deviation of fold RMSE, in volts:

\begin{center}
\begin{tabular}{rrrrrr}
\toprule
Hidden & P/Q mean & P/Q std. & +Vref mean & +Vref std. & Joint \\
\midrule
155 & 0.584702 & 0.002578 & 0.054287 & 0.000410 & 0.319494 \\
93 & 0.593790 & 0.015297 & 0.063779 & 0.003495 & 0.328785 \\
\bottomrule
\end{tabular}
\end{center}

Selected hidden-neuron count: 155.

\subsection*{E2.1: Seed comparison}

\begin{center}
\begin{tabular}{rrrr}
\toprule
Seed & P/Q mean RMSE (V) & +Vref mean RMSE (V) & Joint (V) \\
\midrule
1 & 0.597398 & 0.052285 & 0.324841 \\
0 & 0.600967 & 0.052752 & 0.326860 \\
2 & 0.607313 & 0.054578 & 0.330946 \\
\bottomrule
\end{tabular}
\end{center}

Selected seed: 1.

\subsection*{E2.2: Final training}

\begin{center}
\begin{tabular}{lrrr}
\toprule
Model & Inputs & Outputs & Completed epochs \\
\midrule
P/Q only & 62 & 31 & 40 \\
P/Q + Vref & 65 & 31 & 40 \\
\bottomrule
\end{tabular}
\end{center}

\subsection*{E2.3: Held-out test metrics}

\begin{center}
\begin{tabular}{lrr}
\toprule
Model & RMSE (V) & MAE (V) \\
\midrule
P/Q only & 0.658955 & 0.479735 \\
P/Q + Vref & 0.047638 & 0.034860 \\
\bottomrule
\end{tabular}
\end{center}

\begin{longtable}{rrr}
\caption{Test RMSE for each of the 31 target customers.}\\
\toprule Customer & P/Q only RMSE (V) & P/Q + Vref RMSE (V) \\ \midrule
\endfirsthead
\toprule Customer & P/Q only RMSE (V) & P/Q + Vref RMSE (V) \\ \midrule
\endhead
\bottomrule\endfoot
1 & 0.564115 & 0.018393 \\
2 & 0.709498 & 0.025437 \\
3 & 0.648981 & 0.025153 \\
4 & 0.572872 & 0.023709 \\
5 & 0.752155 & 0.032014 \\
6 & 0.660002 & 0.026363 \\
7 & 0.564298 & 0.027801 \\
8 & 0.736837 & 0.042775 \\
9 & 0.652373 & 0.034059 \\
10 & 0.576141 & 0.029734 \\
11 & 0.739425 & 0.041125 \\
12 & 0.652856 & 0.039666 \\
13 & 0.567449 & 0.037427 \\
14 & 0.746576 & 0.045619 \\
15 & 0.648879 & 0.036716 \\
16 & 0.574454 & 0.044865 \\
17 & 0.735714 & 0.058042 \\
18 & 0.654817 & 0.049783 \\
19 & 0.588425 & 0.055577 \\
20 & 0.744757 & 0.054893 \\
21 & 0.662171 & 0.059403 \\
22 & 0.579683 & 0.057289 \\
23 & 0.750198 & 0.060591 \\
24 & 0.670091 & 0.066627 \\
25 & 0.584982 & 0.054939 \\
26 & 0.751607 & 0.070531 \\
27 & 0.665785 & 0.054701 \\
28 & 0.563807 & 0.056967 \\
29 & 0.745268 & 0.069582 \\
30 & 0.669573 & 0.053266 \\
31 & 0.580494 & 0.054017 \\
\end{longtable}

\subsection*{E2.4: Worst P/Q-only customer}

Customer 5: P/Q-only RMSE 0.752155 V; P/Q + Vref RMSE 0.032014 V.

Running E2.4 displays both requested panels and saves the plot here:
\begin{center}
\path{results/exercises/exercise_worst_customer.png}
\end{center}

These are results under the current 31-target, supplied-reference convention. They do not establish the undocumented physical origin of the reference signals.

\end{document}
